\problema{Shortest Path in Binary Matrix}{1091}{src/01-grafos/1091-shortest-path-binary-matrix.cpp}{M}

En una cuadrícula $n\times n$ de \texttt{0} y \texttt{1}, un \emph{camino
claro} va de la esquina $(0,0)$ a la esquina $(n{-}1,n{-}1)$ pisando solo
celdas con \texttt{0} y moviéndose en las 8 direcciones (incluidas las
diagonales). Su longitud es el número de celdas visitadas. Devolvemos la
longitud del camino más corto, o $-1$ si no existe.

\paragraph{La idea, con un ejemplo de la vida real.} Queremos ir de una esquina
de la cuadrícula a la esquina opuesta pisando solo casillas libres (\texttt{0}), y
usando el camino más corto posible. Aquí sí podemos movernos en diagonal, así
que cada casilla tiene hasta 8 vecinas: es el movimiento del rey en el ajedrez,
que avanza a cualquiera de las casillas que lo rodean.

Para encontrar el camino más corto usamos \emph{BFS}:
soltamos una onda desde la casilla de salida, como los círculos que hace una
piedra al caer al agua. La onda crece de un paso en un paso y toca antes las
casillas cercanas que las lejanas. Por eso, la \emph{primera} vez que la onda
toca la casilla de llegada, lo hace por el camino más corto: no hay forma de
llegar antes. Ese número de casillas recorridas es la respuesta.

\begin{center}
\ttfamily
\begin{tabular}{ccc}
(r-1,c-1) & (r-1,c) & (r-1,c+1)\\
(r,c-1) & (r,c) & (r,c+1)\\
(r+1,c-1) & (r+1,c) & (r+1,c+1)\\
\end{tabular}
\end{center}

\noindent Las 8 celdas que rodean a $(r,c)$. La primera vez que la onda toca
$(n{-}1,n{-}1)$, esa es la distancia mínima en celdas.

\paragraph{Cómo lo resuelve el código, paso a paso.} Usamos una \emph{cola}:
una fila de espera donde el primero que entra es el primero que sale.
\begin{enumerate}
  \item Si la esquina de salida o la de llegada están bloqueadas
        (valen \texttt{1}), devolvemos $-1$ de inmediato: es imposible.
  \item Metemos la esquina de salida a la cola con longitud 1 (ella misma ya
        cuenta como una casilla del camino).
  \item Procesamos la cola por capas. Cada capa es un paso más de la onda.
        Sacamos una casilla y miramos sus 8 vecinas; si una vecina está libre y
        no la hemos visitado, le anotamos la longitud del camino (la actual más
        1) y la metemos a la cola.
  \item Para no volver a pisar una casilla ya visitada, la marcamos cambiándola
        a \texttt{1} en cuanto la usamos. En cuanto sacamos la esquina de
        llegada de la cola, devolvemos su longitud.
  \item Si vaciamos la cola sin haber llegado nunca, devolvemos $-1$.
\end{enumerate}

\lstinputlisting{src/01-grafos/1091-shortest-path-binary-matrix.cpp}

\complejidad{$O(n^2)$}{$O(n^2)$}

\paragraph{¿Qué significa esa complejidad?} La cuadrícula es de $n\times n$, así
que tiene $n^2$ casillas en total. Que el tiempo sea $O(n^2)$ significa que, en
el peor caso, cada casilla entra y sale de la cola una sola vez: revisamos cada
una una vez. No se puede menos, porque tal vez haya que recorrer todo el
cuadrícula para encontrar la salida.

\paragraph{Consejos para la entrevista (Amazon).} Explica por qué usamos BFS y
no DFS: el DFS encuentra \emph{algún} camino, pero no garantiza que sea el más
corto; el BFS
sí, porque avanza por capas parejas. Cuida los casos raros: que la salida o la
llegada estén bloqueadas (respuesta $-1$ de inmediato) y una cuadrícula de una sola
casilla con un \texttt{0} (respuesta 1). Si además pusieran «costos» distintos
por casilla, el siguiente paso sería usar el algoritmo de Dijkstra o el 0-1
BFS.
